% --- Archon physics formalization source begin ---
% archon:physics
% archon:covers PhyXMiniProblems/problem_phyx_mini_0262.lean
% archon:source-report reports/phyx_mini/problem_phyx_mini_0262.source.json

\chapter{Physics problem phyx\_mini\_0262}
\label{ch:PhyXMiniProblems_problem_phyx_mini_0262}

\paragraph{Problem source.}
\#\# Physical scenario

In figure, the pendulum consists of a uniform disk with radius \$r = 10.0\textbackslash{} cm\$ and mass \$500\textbackslash{} g\$ attached to a uniform rod with length \$L = 500\textbackslash{} mm\$ and mass \$270\textbackslash{} g\$.

\#\# Question

Calculate the period of oscillation.

\#\# Answer choices

- A: A: 1.0 s

- B: B: 1.2 s

- C: C: 1.5 s

- D: D: 1.8 s

\#\# Figure caption (auxiliary; use the image as primary evidence)

The image illustrates a mechanical setup involving a pendulum-like structure. Key components include:

1. **Pivot Point**: At the top, attached to a fixed support (in beige), the pendulum is connected, allowing for rotational movement.

2. **Rod (Pendulum Arm)**: A blue, cylindrical rod extends downward from the pivot. The length of this rod is marked as \textbackslash{}( L \textbackslash{}), which is labeled to indicate its measurement from the pivot to another point further down.

3. **Wheel or Disc**: Attached at the lower end of the rod is a red wheel or disc, shaded to imply depth or a three-dimensional appearance. It's depicted with a central axis of rotation.

4. **Distance \textbackslash{}( r \textbackslash{})**: A labeled line \textbackslash{}( r \textbackslash{}) extends from the center of the wheel to a specific point on its circumference. This indicates the radius of the wheel.

The figure seems to illustrate a pendulum or wheel system, possibly for educational or analytical purposes involving rotational dynamics or mechanics.

\#\# Dataset metadata

Resonance and Harmonics; Physical Model Grounding Reasoning, Multi-Formula Reasoning

\paragraph{Recorded answer/context.}
C: C: 1.5 s

\paragraph{Figure/image path.}
/root/proposal\_for\_physic/science-mango/phyx\_mini\_run/phyx\_data/test\_image/262.png

\paragraph{Formalization target.}
create a compiling Lean file with sorry bodies at `PhyXMiniProblems/problem\_phyx\_mini\_0262.lean`. The Lean declarations must preserve the physical quantities, dimensions or dimensional roles, figure labels, governing-law hypotheses, and final relation expressed by this problem.
Use Mathlib/Physlib names found through LeanExplore where available. If a physics API is missing, introduce faithful local abstractions rather than scalar placeholder aliases.

\begin{theorem}[Physics formalization target]
\label{thm:physics:phyx_mini_0262:target}
\lean{PhyXMiniProblems.ProblemPhyXMini0262.rodDiskLinearizedPeriod_matches_recordedAnswerC}
\uses{def:physics:phyx-mini-0262:phyxminiproblems-problemphyxmini0262-massinkilograms, def:physics:phyx-mini-0262:phyxminiproblems-problemphyxmini0262-lengthinmeters, def:physics:phyx-mini-0262:phyxminiproblems-problemphyxmini0262-timeinseconds, def:physics:phyx-mini-0262:phyxminiproblems-problemphyxmini0262-accelerationinmeterspersecondsquared, def:physics:phyx-mini-0262:phyxminiproblems-problemphyxmini0262-roddiskpendulumsetup, def:physics:phyx-mini-0262:phyxminiproblems-problemphyxmini0262-matchesprimaryroddiskfigure, def:physics:phyx-mini-0262:phyxminiproblems-problemphyxmini0262-hasroddiskpendulumproblemdata, def:physics:phyx-mini-0262:phyxminiproblems-problemphyxmini0262-satisfiesroddiskcompoundpendulumlaws, def:physics:phyx-mini-0262:phyxminiproblems-problemphyxmini0262-answerchoice, def:physics:phyx-mini-0262:phyxminiproblems-problemphyxmini0262-recordedanswerchoice, def:physics:phyx-mini-0262:phyxminiproblems-problemphyxmini0262-matchesdisplayedperiod}
The assigned autoformalize agent should translate this physics problem into Lean declarations in the covered file, with theorem and lemma proof bodies written as `by sorry`.
\end{theorem}
\begin{proof}
This is an autoformalization task, not a proof task. Produce statements that can later be proved by the physics prover without weakening the physical model.
\end{proof}

% --- Archon declaration topology begin ---
\section{Declaration topology}
\begin{definition}[Mass Quantity]
  \label{def:physics:phyx-mini-0262:phyxminiproblems-problemphyxmini0262-massquantity}
  \lean{PhyXMiniProblems.ProblemPhyXMini0262.MassQuantity}
  A nonnegative physical mass
\end{definition}
\begin{proof}
  This is the declared model definition of Mass Quantity.
\end{proof}
\begin{definition}[Length Quantity]
  \label{def:physics:phyx-mini-0262:phyxminiproblems-problemphyxmini0262-lengthquantity}
  \lean{PhyXMiniProblems.ProblemPhyXMini0262.LengthQuantity}
  A nonnegative physical length
\end{definition}
\begin{proof}
  This is the declared model definition of Length Quantity.
\end{proof}
\begin{definition}[Time Quantity]
  \label{def:physics:phyx-mini-0262:phyxminiproblems-problemphyxmini0262-timequantity}
  \lean{PhyXMiniProblems.ProblemPhyXMini0262.TimeQuantity}
  A nonnegative physical duration
\end{definition}
\begin{proof}
  This is the declared model definition of Time Quantity.
\end{proof}
\begin{definition}[Acceleration Magnitude Quantity]
  \label{def:physics:phyx-mini-0262:phyxminiproblems-problemphyxmini0262-accelerationmagnitudequantity}
  \lean{PhyXMiniProblems.ProblemPhyXMini0262.AccelerationMagnitudeQuantity}
  A nonnegative acceleration magnitude, with dimension L T⁻²
\end{definition}
\begin{proof}
  This is the declared model definition of Acceleration Magnitude Quantity.
\end{proof}
\begin{definition}[Moment Of Inertia Quantity]
  \label{def:physics:phyx-mini-0262:phyxminiproblems-problemphyxmini0262-momentofinertiaquantity}
  \lean{PhyXMiniProblems.ProblemPhyXMini0262.MomentOfInertiaQuantity}
  A moment of inertia about an axis, with dimension M L²
\end{definition}
\begin{proof}
  This is the declared model definition of Moment Of Inertia Quantity.
\end{proof}
\begin{definition}[Restoring Torque Coefficient Quantity]
  \label{def:physics:phyx-mini-0262:phyxminiproblems-problemphyxmini0262-restoringtorquecoefficientquantity}
  \lean{PhyXMiniProblems.ProblemPhyXMini0262.RestoringTorqueCoefficientQuantity}
  The nonnegative coefficient κ in the exact gravitational restoring torque τ(θ) = -κ sin θ. Radian measure is dimensionless, so κ has dimension M L² T⁻²
\end{definition}
\begin{proof}
  This is the declared model definition of Restoring Torque Coefficient Quantity.
\end{proof}
\begin{definition}[Torque Quantity]
  \label{def:physics:phyx-mini-0262:phyxminiproblems-problemphyxmini0262-torquequantity}
  \lean{PhyXMiniProblems.ProblemPhyXMini0262.TorqueQuantity}
  A signed torque, with dimension M L² T⁻²
\end{definition}
\begin{proof}
  This is the declared model definition of Torque Quantity.
\end{proof}
\begin{definition}[mass In Kilograms]
  \label{def:physics:phyx-mini-0262:phyxminiproblems-problemphyxmini0262-massinkilograms}
  \lean{PhyXMiniProblems.ProblemPhyXMini0262.massInKilograms}
  \uses{def:physics:phyx-mini-0262:phyxminiproblems-problemphyxmini0262-massquantity}
  Read a physical mass in kilograms
\end{definition}
\begin{proof}
  This is the declared model definition of mass In Kilograms.
\end{proof}
\begin{definition}[length In Meters]
  \label{def:physics:phyx-mini-0262:phyxminiproblems-problemphyxmini0262-lengthinmeters}
  \lean{PhyXMiniProblems.ProblemPhyXMini0262.lengthInMeters}
  \uses{def:physics:phyx-mini-0262:phyxminiproblems-problemphyxmini0262-lengthquantity}
  Read a physical length in meters
\end{definition}
\begin{proof}
  This is the declared model definition of length In Meters.
\end{proof}
\begin{definition}[time In Seconds]
  \label{def:physics:phyx-mini-0262:phyxminiproblems-problemphyxmini0262-timeinseconds}
  \lean{PhyXMiniProblems.ProblemPhyXMini0262.timeInSeconds}
  \uses{def:physics:phyx-mini-0262:phyxminiproblems-problemphyxmini0262-timequantity}
  Read a physical duration in seconds
\end{definition}
\begin{proof}
  This is the declared model definition of time In Seconds.
\end{proof}
\begin{definition}[acceleration In Meters Per Second Squared]
  \label{def:physics:phyx-mini-0262:phyxminiproblems-problemphyxmini0262-accelerationinmeterspersecondsquared}
  \lean{PhyXMiniProblems.ProblemPhyXMini0262.accelerationInMetersPerSecondSquared}
  \uses{def:physics:phyx-mini-0262:phyxminiproblems-problemphyxmini0262-accelerationmagnitudequantity}
  Read an acceleration magnitude in meters per second squared
\end{definition}
\begin{proof}
  This is the declared model definition of acceleration In Meters Per Second Squared.
\end{proof}
\begin{definition}[moment Of Inertia In Kilogram Meters Squared]
  \label{def:physics:phyx-mini-0262:phyxminiproblems-problemphyxmini0262-momentofinertiainkilogrammeterssquared}
  \lean{PhyXMiniProblems.ProblemPhyXMini0262.momentOfInertiaInKilogramMetersSquared}
  \uses{def:physics:phyx-mini-0262:phyxminiproblems-problemphyxmini0262-momentofinertiaquantity}
  Read a moment of inertia in kilogram meter squared
\end{definition}
\begin{proof}
  This is the declared model definition of moment Of Inertia In Kilogram Meters Squared.
\end{proof}
\begin{definition}[restoring Coefficient In Newton Meters]
  \label{def:physics:phyx-mini-0262:phyxminiproblems-problemphyxmini0262-restoringcoefficientinnewtonmeters}
  \lean{PhyXMiniProblems.ProblemPhyXMini0262.restoringCoefficientInNewtonMeters}
  \uses{def:physics:phyx-mini-0262:phyxminiproblems-problemphyxmini0262-restoringtorquecoefficientquantity}
  Read a restoring-torque coefficient in newton meters
\end{definition}
\begin{proof}
  This is the declared model definition of restoring Coefficient In Newton Meters.
\end{proof}
\begin{definition}[torque In Newton Meters]
  \label{def:physics:phyx-mini-0262:phyxminiproblems-problemphyxmini0262-torqueinnewtonmeters}
  \lean{PhyXMiniProblems.ProblemPhyXMini0262.torqueInNewtonMeters}
  \uses{def:physics:phyx-mini-0262:phyxminiproblems-problemphyxmini0262-torquequantity}
  Read a signed torque in newton meters
\end{definition}
\begin{proof}
  This is the declared model definition of torque In Newton Meters.
\end{proof}
\begin{definition}[Pendulum Component]
  \label{def:physics:phyx-mini-0262:phyxminiproblems-problemphyxmini0262-pendulumcomponent}
  \lean{PhyXMiniProblems.ProblemPhyXMini0262.PendulumComponent}
  The two massive components of the compound pendulum
\end{definition}
\begin{proof}
  This is the declared model definition of Pendulum Component.
\end{proof}
\begin{definition}[Figure Point]
  \label{def:physics:phyx-mini-0262:phyxminiproblems-problemphyxmini0262-figurepoint}
  \lean{PhyXMiniProblems.ProblemPhyXMini0262.FigurePoint}
  Distinguished locations used by the dimension arrows in the image
\end{definition}
\begin{proof}
  This is the declared model definition of Figure Point.
\end{proof}
\begin{definition}[Figure Label]
  \label{def:physics:phyx-mini-0262:phyxminiproblems-problemphyxmini0262-figurelabel}
  \lean{PhyXMiniProblems.ProblemPhyXMini0262.FigureLabel}
  The two mathematical labels printed in the primary image
\end{definition}
\begin{proof}
  This is the declared model definition of Figure Label.
\end{proof}
\begin{definition}[Rod Mass Distribution]
  \label{def:physics:phyx-mini-0262:phyxminiproblems-problemphyxmini0262-rodmassdistribution}
  \lean{PhyXMiniProblems.ProblemPhyXMini0262.RodMassDistribution}
  Idealized mass distribution of the rod
\end{definition}
\begin{proof}
  This is the declared model definition of Rod Mass Distribution.
\end{proof}
\begin{definition}[Disk Mass Distribution]
  \label{def:physics:phyx-mini-0262:phyxminiproblems-problemphyxmini0262-diskmassdistribution}
  \lean{PhyXMiniProblems.ProblemPhyXMini0262.DiskMassDistribution}
  Idealized mass distribution of the wheel-like body
\end{definition}
\begin{proof}
  This is the declared model definition of Disk Mass Distribution.
\end{proof}
\begin{definition}[Rod Disk Attachment]
  \label{def:physics:phyx-mini-0262:phyxminiproblems-problemphyxmini0262-roddiskattachment}
  \lean{PhyXMiniProblems.ProblemPhyXMini0262.RodDiskAttachment}
  How the lower end of the rod is joined to the disk
\end{definition}
\begin{proof}
  This is the declared model definition of Rod Disk Attachment.
\end{proof}
\begin{definition}[Compound Pendulum Figure]
  \label{def:physics:phyx-mini-0262:phyxminiproblems-problemphyxmini0262-compoundpendulumfigure}
  \lean{PhyXMiniProblems.ProblemPhyXMini0262.CompoundPendulumFigure}
  \uses{def:physics:phyx-mini-0262:phyxminiproblems-problemphyxmini0262-figurepoint, def:physics:phyx-mini-0262:phyxminiproblems-problemphyxmini0262-figurelabel}
  Qualitative evidence and dimension endpoints visible in the image
\end{definition}
\begin{proof}
  This is the declared model definition of Compound Pendulum Figure.
\end{proof}
\begin{definition}[Rod Disk Pendulum Setup]
  \label{def:physics:phyx-mini-0262:phyxminiproblems-problemphyxmini0262-roddiskpendulumsetup}
  \lean{PhyXMiniProblems.ProblemPhyXMini0262.RodDiskPendulumSetup}
  \uses{def:physics:phyx-mini-0262:phyxminiproblems-problemphyxmini0262-massquantity, def:physics:phyx-mini-0262:phyxminiproblems-problemphyxmini0262-lengthquantity, def:physics:phyx-mini-0262:phyxminiproblems-problemphyxmini0262-timequantity, def:physics:phyx-mini-0262:phyxminiproblems-problemphyxmini0262-accelerationmagnitudequantity, def:physics:phyx-mini-0262:phyxminiproblems-problemphyxmini0262-momentofinertiaquantity, def:physics:phyx-mini-0262:phyxminiproblems-problemphyxmini0262-restoringtorquecoefficientquantity, def:physics:phyx-mini-0262:phyxminiproblems-problemphyxmini0262-torquequantity, def:physics:phyx-mini-0262:phyxminiproblems-problemphyxmini0262-pendulumcomponent, def:physics:phyx-mini-0262:phyxminiproblems-problemphyxmini0262-rodmassdistribution, def:physics:phyx-mini-0262:phyxminiproblems-problemphyxmini0262-diskmassdistribution, def:physics:phyx-mini-0262:phyxminiproblems-problemphyxmini0262-roddiskattachment, def:physics:phyx-mini-0262:phyxminiproblems-problemphyxmini0262-compoundpendulumfigure}
  Physical quantities for the rod--disk compound pendulum. The nonlinear period function and its tangent-linearized zero-amplitude limit are distinct fields. Neither is defined using a displayed answer
\end{definition}
\begin{proof}
  This is the declared model definition of Rod Disk Pendulum Setup.
\end{proof}
\begin{definition}[Matches Primary Rod Disk Figure]
  \label{def:physics:phyx-mini-0262:phyxminiproblems-problemphyxmini0262-matchesprimaryroddiskfigure}
  \lean{PhyXMiniProblems.ProblemPhyXMini0262.MatchesPrimaryRodDiskFigure}
  \uses{def:physics:phyx-mini-0262:phyxminiproblems-problemphyxmini0262-lengthinmeters, def:physics:phyx-mini-0262:phyxminiproblems-problemphyxmini0262-roddiskpendulumsetup}
  Primary-image readout. The L arrow runs from the fixed pivot to the rod--disk rim attachment, while the r arrow continues from that rim point to the disk center. The resulting lever arm L + r is geometry, not the requested period
\end{definition}
\begin{proof}
  This is the declared model definition of Matches Primary Rod Disk Figure.
\end{proof}
\begin{definition}[Has Rod Disk Pendulum Problem Data]
  \label{def:physics:phyx-mini-0262:phyxminiproblems-problemphyxmini0262-hasroddiskpendulumproblemdata}
  \lean{PhyXMiniProblems.ProblemPhyXMini0262.HasRodDiskPendulumProblemData}
  \uses{def:physics:phyx-mini-0262:phyxminiproblems-problemphyxmini0262-massinkilograms, def:physics:phyx-mini-0262:phyxminiproblems-problemphyxmini0262-lengthinmeters, def:physics:phyx-mini-0262:phyxminiproblems-problemphyxmini0262-accelerationinmeterspersecondsquared, def:physics:phyx-mini-0262:phyxminiproblems-problemphyxmini0262-roddiskpendulumsetup}
  Numerical and qualitative data supplied by the problem. Centimeters, millimeters, and grams are converted to coherent SI readouts. The standard value 9.81 m/s² is the environmental datum used to compare the displayed choices
\end{definition}
\begin{proof}
  This is the declared model definition of Has Rod Disk Pendulum Problem Data.
\end{proof}
\begin{definition}[Satisfies Rod Disk Compound Pendulum Laws]
  \label{def:physics:phyx-mini-0262:phyxminiproblems-problemphyxmini0262-satisfiesroddiskcompoundpendulumlaws}
  \lean{PhyXMiniProblems.ProblemPhyXMini0262.SatisfiesRodDiskCompoundPendulumLaws}
  \uses{def:physics:phyx-mini-0262:phyxminiproblems-problemphyxmini0262-massinkilograms, def:physics:phyx-mini-0262:phyxminiproblems-problemphyxmini0262-lengthinmeters, def:physics:phyx-mini-0262:phyxminiproblems-problemphyxmini0262-timeinseconds, def:physics:phyx-mini-0262:phyxminiproblems-problemphyxmini0262-accelerationinmeterspersecondsquared, def:physics:phyx-mini-0262:phyxminiproblems-problemphyxmini0262-momentofinertiainkilogrammeterssquared, def:physics:phyx-mini-0262:phyxminiproblems-problemphyxmini0262-restoringcoefficientinnewtonmeters, def:physics:phyx-mini-0262:phyxminiproblems-problemphyxmini0262-torqueinnewtonmeters, def:physics:phyx-mini-0262:phyxminiproblems-problemphyxmini0262-roddiskpendulumsetup}
  Rigid-body, gravitational, and local small-oscillation relations. The rod has pivot inertia m L² / 3. The disk has central inertia m r² / 2, and the parallel-axis theorem translates it through d = L + r. The exact signed gravitational torque is -κ sin θ. The HasDerivAt and little-o fields make the replacement of that torque by -κ θ explicitly local at the stable equilibrium. The frequency and period laws refer only to the tangent-linearized oscillator.
\end{definition}
\begin{proof}
  This is the declared model definition of Satisfies Rod Disk Compound Pendulum Laws.
\end{proof}
\begin{definition}[Answer Choice]
  \label{def:physics:phyx-mini-0262:phyxminiproblems-problemphyxmini0262-answerchoice}
  \lean{PhyXMiniProblems.ProblemPhyXMini0262.AnswerChoice}
  Labels of the four displayed period choices
\end{definition}
\begin{proof}
  This is the declared model definition of Answer Choice.
\end{proof}
\begin{definition}[seconds]
  \label{def:physics:phyx-mini-0262:phyxminiproblems-problemphyxmini0262-answerchoice-seconds}
  \lean{PhyXMiniProblems.ProblemPhyXMini0262.AnswerChoice.seconds}
  \uses{def:physics:phyx-mini-0262:phyxminiproblems-problemphyxmini0262-answerchoice}
  The period in seconds printed beside each answer label
\end{definition}
\begin{proof}
  This is the declared model definition of seconds.
\end{proof}
\begin{definition}[recorded Answer Choice]
  \label{def:physics:phyx-mini-0262:phyxminiproblems-problemphyxmini0262-recordedanswerchoice}
  \lean{PhyXMiniProblems.ProblemPhyXMini0262.recordedAnswerChoice}
  \uses{def:physics:phyx-mini-0262:phyxminiproblems-problemphyxmini0262-answerchoice}
  The answer label recorded in the supplied dataset
\end{definition}
\begin{proof}
  This is the declared model definition of recorded Answer Choice.
\end{proof}
\begin{definition}[Matches Displayed Period]
  \label{def:physics:phyx-mini-0262:phyxminiproblems-problemphyxmini0262-matchesdisplayedperiod}
  \lean{PhyXMiniProblems.ProblemPhyXMini0262.MatchesDisplayedPeriod}
  \uses{def:physics:phyx-mini-0262:phyxminiproblems-problemphyxmini0262-timequantity, def:physics:phyx-mini-0262:phyxminiproblems-problemphyxmini0262-timeinseconds, def:physics:phyx-mini-0262:phyxminiproblems-problemphyxmini0262-answerchoice}
  Agreement with an answer displayed to the nearest tenth of a second
\end{definition}
\begin{proof}
  This is the declared model definition of Matches Displayed Period.
\end{proof}
% --- Archon declaration topology end ---
% --- Archon physics formalization source end ---
